<<<<<<< HEAD
\documentclass[10pt, spanish]{beamer}

% --- PAQUETES ---
\usetheme{metropolis}           
\usepackage[utf8]{inputenc}
\usepackage[spanish]{babel}
\usepackage{listings}
\usepackage{xcolor}
\usepackage{graphicx}
\usepackage{amsmath, amssymb}
\usepackage{karnaugh-map} % Para mapas de Karnaugh
\usepackage{tikz}         % Para diagramas y circuitos
\usepackage{circuitikz}   % Para circuitos eléctricos
\usepackage{booktabs}     % Para tablas elegantes
\usepackage{hyperref}
\usepackage[backend=biber]{biblatex} % Configura Biber como motor
\addbibresource{bibliography.bib}    % Nombre exacto de tu archivo .bib

% --- CONFIGURACIÓN DE CÓDIGO (LISTINGS) ---
\definecolor{codegreen}{rgb}{0,0.6,0}
\definecolor{codegray}{rgb}{0.5,0.5,0.5}
\definecolor{backcolour}{rgb}{0.95,0.95,0.92}

\lstset{
    backgroundcolor=\color{backcolour},   
    commentstyle=\color{codegreen},
    keywordstyle=\color{blue},
    numberstyle=\tiny\color{codegray},
    stringstyle=\color{red},
    basicstyle=\ttfamily\tiny,
    breaklines=true,
    language=[LaTeX]TeX,
    frame=single
}

\title{\LaTeX{} para Ingeniería y Academia}
\subtitle{De la Estructura APA/IEEE}
\author{Introducción a la Composición Tipográfica Profesional}
\date{\today}

\begin{document}

\maketitle

\section{Introducción}

\begin{frame}{¿Qué es \LaTeX? y ¿Por qué usarlo?}
    \begin{itemize}
        \item \textbf{Concepto:} No es un procesador de texto (como Word), es un lenguaje de marcado para composición tipográfica (WYSIWYM).
        \item \textbf{Ventajas:}
            \begin{itemize}
                \item Manejo automático de bibliografía e índices.
                \item Calidad matemática insuperable.
                \item Estabilidad en documentos largos (tesis de 300+ páginas).
                \item Separación total de \textbf{Contenido} y \textbf{Formato}.
            \end{itemize}
    \end{itemize}
\end{frame}

\section{Formatos Estándar: APA 7 e IEEE}

\begin{frame}[fragile]{Estructura APA 7ma Edición}
    Se usa el paquete \texttt{apa7}. Sus modos principales son:
    \begin{itemize}
        \item \textbf{man}: (manuscript) Doble espacio, una columna (para revisión).
        \item \textbf{stu}: (student) Formato para trabajos estudiantiles.
        \item \textbf{jou}: (journal) Formato de revista profesional.
    \end{itemize}
\begin{lstlisting}
\documentclass[stu, spanish]{apa7}
\title{Analisis de Circuitos}
\shorttitle{APA7 en Beamer}
\authorsname{Juan Perez}
\affiliation{Universidad X}
\begin{document}
    \maketitle
    \section{Introduccion} ...
\end{document}
\end{lstlisting}
\end{frame}

\begin{frame}[fragile]{Estructura IEEE}
    Ideal para artículos de ingeniería. Clases comunes:
    \begin{itemize}
        \item \textbf{conference}: Para ponencias en congresos.
        \item \textbf{journal}: Para publicaciones en revistas IEEE.
        \item \textbf{comsoc}: Para comunicaciones.
    \end{itemize}
\begin{lstlisting}
\documentclass[conference]{IEEEtran}
\title{Implementacion de Logica Digital}
\author{\IEEEauthorblockN{Maria Garcia} \IEEEauthorblockA{Escuela de Ingenieria}}
\begin{document}
    \maketitle
    \begin{abstract} Resumen del articulo... \end{abstract}
\end{document}
\end{lstlisting}
\end{frame}

\section{Elementos Académicos Comunes}

\begin{frame}[fragile]{Ecuaciones y Referencias}
    Diferencias fundamentales en matemáticas:
    \begin{itemize}
        \item \textbf{En línea:} \lstinline|$E=mc^2$| para el texto ($E=mc^2$).
        \item \textbf{Bloque:} \lstinline|$$...$$| o \lstinline|\[...\]| no se numeran.
        \item \textbf{Entorno \texttt{equation}:} Se numera y permite usar \texttt{\textbackslash label}.
    \end{itemize}
\begin{lstlisting}
\begin{equation}
    V = I \cdot R \label{eq:ohm}
\end{equation}
Como se observa en la ec. (\ref{eq:ohm})...
\end{lstlisting}
    \textbf{Salida:}
    \begin{equation}
        V = I \cdot R \label{eq:ohm_example}
    \end{equation}
\end{frame}

\begin{frame}[fragile]{Símbolos y Letras Griegas}
    Para símbolos especiales, visita: \href{https://manualdelatex.com/simbolos#chapter3}{Manual de Símbolos \LaTeX}.
    \begin{itemize}
        \item Letras: \lstinline|$\alpha, \beta, \gamma, \Delta, \Omega$| $\rightarrow$ $\alpha, \beta, \gamma, \Delta, \Omega$.
        \item Operadores: \lstinline|$\int, \sum, \infty, \nabla$| $\rightarrow$ $\int, \sum, \infty, \nabla$.
    \end{itemize}
    Para fórmulas complejas sin escribir código: \href{https://editor-codecogs-com.translate.goog/?_x_tr_sl=en&_x_tr_tl=es&_x_tr_hl=es&_x_tr_pto=tc}{Editor Visual CodeCogs}.
\end{frame}

\begin{frame}[fragile]{Tablas de Cálculo de Errores}
    Recomiendo usar: \href{https://www.tablesgenerator.com/}{Tables Generator}.
\begin{lstlisting}
\begin{table}[h]
\centering
\caption{Comparativa de tensiones.}
\label{tab:errores}
\begin{tabular}{|c|c|c|c|}
\hline
\textbf{Nodo} & \textbf{Medido (V)} & \textbf{Simulado (V)} & \textbf{Error (\%)} \\ \hline
1 & 4.98 & 5.00 & 0.4\% \\ \hline
2 & 2.45 & 2.50 & 2.0\% \\ \hline
\end{tabular}
\end{table}
\end{lstlisting}

También se pueden apoyar de la IA para generar el código de tablas más complejas.
\end{frame}

\begin{frame}[fragile]{Inserción de Imágenes y Figuras}
Para incluir imágenes se usa el paquete \texttt{graphicx} y el entorno \texttt{figure}.
\begin{lstlisting}
\begin{figure}[htbp] % Opciones de posicionamiento
    \caption{Diagrama del sistema propuesto.} % Caption ARRIBA
    \centering
    \includegraphics[width=0.5\textwidth]{mi_imagen.png}
    \label{fig:diagrama} % Etiqueta para referenciar
    \par\vspace{2pt}\small \textit{Nota: Datos obtenidos en laboratorio.} % Nota abajo
\end{figure}
Como se aprecia en la Fig. \ref{fig:diagrama}, los resultados...
\end{lstlisting}

\textbf{Opciones de Posicionamiento [!htbp]:}
\begin{itemize}
    \item \textbf{h} (here): Intenta ponerla justo donde está el código.
    \item \textbf{t} (top): Al principio de la página actual.
    \item \textbf{b} (bottom): Al final de la página actual.
    \item \textbf{p} (page): En una página dedicada solo a flotantes.
    \item \textbf{!} : Fuerza a \LaTeX{} a ignorar algunas reglas estéticas.
\end{itemize}
\end{frame}

\begin{frame}[fragile]{Citas en el Texto: APA 7 vs. IEEE}
La forma de citar varía drásticamente según el estándar:

\begin{table}[]
\tiny
\begin{tabular}{@{}lll@{}}
\toprule
\textbf{Tipo de Cita} & \textbf{Normas APA 7ma Ed.} & \textbf{Normas IEEE} \\ \midrule
\textbf{Parentética} & (García \& Pérez, 2023) & [1] \\
\textbf{Narrativa} & García y Pérez (2023) afirman... & García y Pérez en [1]... \\
\textbf{Múltiple} & (Altamirano, 2020; Zappa, 2018) & [2], [5] \\
\textbf{Páginas} & (García, 2023, p. 45) & [1, p. 45] \\ \bottomrule
\end{tabular}
\end{table}

\textbf{Comandos en \LaTeX:}
\begin{itemize}
    \item \lstinline|\cite{key}|: Cita estándar ([1] o (García, 2023)).
    \item \lstinline|\textcite{key}|: Cita narrativa (García (2023)). Solo en APA/BibLaTeX.
    \item \lstinline|\citeonline{key}|: Variante para IEEE según algunos paquetes.
\end{itemize}
\end{frame}


\begin{frame}[fragile]{Comandos de Citación: \texttt{\textbackslash parentcite} vs \texttt{\textbackslash textcite}}
Para documentos en \textbf{APA 7}, no basta con usar \texttt{\textbackslash cite}. \texttt{biblatex} ofrece comandos específicos para cumplir con la norma:

\begin{itemize}
    \item \textbf{\textbackslash parentcite\{key\}}: Cita tradicional al final de un párrafo.
    \begin{itemize}
        \item \textit{Resultado APA:} (García \& Pérez, 2024)
        \item \textit{Resultado IEEE:} [1]
    \end{itemize}
    \item \textbf{\textbackslash textcite\{key\}}: Cita narrativa donde el autor es el sujeto.
    \begin{itemize}
        \item \textit{Resultado APA:} García y Pérez (2024)
        \item \textit{Resultado IEEE:} García y Pérez [1]
    \end{itemize}
\end{itemize}

\textbf{Ejemplo de uso en el código:}
\begin{lstlisting}
Segun \textcite{manual_latex}, los simbolos son ...
Esto se debe a la estructura del nucelo \parencite{codecogs}.
\end{lstlisting}

\textit{Nota: Usar estos comandos hace que tu documento sea ``estilo-independiente''; si cambias de APA a IEEE, \LaTeX{} ajustará el formato automáticamente.}
\end{frame}


\begin{frame}[fragile]{Configuración de Bibliografía Externa}
Para gestionar referencias, se recomienda el paquete \texttt{biblatex}. En el \textbf{preámbulo} (antes de \texttt{\textbackslash begin\{document\}}) debes agregar:

\begin{lstlisting}
% 1. Cargar el paquete y definir el estilo (apa, ieee, etc.)
\usepackage[style=apa, backend=biber]{biblatex}

% 2. Vincular el archivo .bib (IMPORTANTE: incluir extension)
\addbibresource{referencias.bib}
\end{lstlisting}

\textbf{Puntos clave:}
\begin{itemize}
    \item \textbf{style:} Define si se verá como (García, 2024) o [1].
    \item \textbf{backend=biber:} Es el motor que procesa los datos (más potente que el antiguo BibTeX).
    \item \textbf{Comando de salida:} Usa \texttt{\textbackslash printbibliography} donde quieras que aparezca la lista final.
\end{itemize}
\end{frame}


\begin{frame}[fragile]{Estructura del archivo \texttt{referencias.bib}}
El archivo bibliográfico es una base de datos en texto plano. Cada entrada tiene un tipo (\texttt{@online}, \texttt{@article}, etc.) y una \textbf{llave} única.

\begin{lstlisting}
@online{manual_latex,
    author = {Manual de LaTeX},
    title  = {Simbolos Matematicos},
    year   = {2024},
    url    = {https://manualdelatex.com/simbolos},
    urldate = {2024-05-20}
}

@online{codecogs,
    author = {CodeCogs},
    title  = {Editor Visual de Ecuaciones},
    year   = {2024},
    url    = {https://editor.codecogs.com}
}

@online{tablesgen,
    author = {Tables Generator},
    title  = {Generador de tablas para LaTeX},
    url    = {https://www.tablesgenerator.com}
}
\end{lstlisting}
\textit{Nota: La "llave" (ej. \texttt{codecogs}) es lo que usas para citar con \texttt{\textbackslash cite\{codecogs\}}.}
\end{frame}



\section{Lógica Digital y Gráficos}

\begin{frame}[fragile]{Mapas de K. con Minterminos y Condiciones ``No Importa''}
Podemos optimizar el llenado con \texttt{\textbackslash minterms}, \texttt{\textbackslash maxterms} y usar \textbf{X} (\texttt{\textbackslash autoterms\{X\}}) para simplificar.

\begin{columns}
\begin{column}{0.5\textwidth}
\begin{lstlisting}
\begin{karnaugh-map}[4][4][1][$CD$][$AB$]
	\minterms{4,10,12,14}
	\maxterms{2,3,6,8}
	\autoterms[X]
	\implicant{4}{12}
	\implicant{14}{10}
\end{karnaugh-map}
\end{lstlisting}
\end{column}
\begin{column}{0.5\textwidth}
    \centering
	\begin{karnaugh-map}[4][4][1][$CD$][$AB$]
    		\minterms{4,10,12,14}
	    \maxterms{2,3,6,8}
    		\autoterms[X]
	    \implicant{4}{12} % Grupo de las 4 esquinas
	    \implicant{14}{10}  % Grupo de 2 (celdas 5 y 7)
	\end{karnaugh-map}
\end{column}
\end{columns}

%\small \textbf{Nota:} \texttt{\textbackslash implicant\{0\}\{10\}} en un mapa de 4x4 identifica que los extremos son adyacentes y crea los arcos en las esquinas.
\end{frame}

\begin{frame}[fragile]{El poder de TikZ y Circuitikz}
    Permite dibujar circuitos directamente en el código.
\begin{lstlisting}
\begin{circuitikz}
  \draw (0,0) to[V=$V_s$] (0,2) to[R=$R_1$] (2,2) 
        to[C=$C_1$] (2,0) -- (0,0);
\end{circuitikz}
\end{lstlisting}
    \centering
    \begin{circuitikz}[scale=0.7]
      \draw (0,0) to[V=$V_s$] (0,2) to[R=$10\Omega$] (2,2) 
            to[C=$1\mu F$] (2,0) -- (0,0);
    \end{circuitikz}
\end{frame}

\section{Instalación y Compilación}

\begin{frame}{Compilación Local y Herramientas}
    Para trabajar sin internet necesitas:
    \begin{enumerate}
        \item \textbf{Un Motor (\LaTeX{} Distribution):}
            \begin{itemize}
                \item \textbf{Windows:} MiKTeX (ligero) o TeX Live.
                \item \textbf{Linux:} TeX Live (\texttt{sudo apt install texlive-full}).
            \end{itemize}
        \item \textbf{Un Editor (IDE):}
            \begin{itemize}
                \item \textbf{Texmaker:} Multiplataforma, muy estable y gratuito.
                \item \textbf{TeXstudio:} Con autocompletado avanzado.
                \item \textbf{VS Code:} Con la extensión "LaTeX Workshop".
            \end{itemize}
    \end{enumerate}
\end{frame}

\begin{frame}[standout]
    ¡Gracias! \\
    ¿Preguntas?
\end{frame}

\section{Referencias Bibliográficas}

\begin{frame}{Referencias en Formato IEEE}
    % Nota: Para mostrar esto realmente, tendrias que cambiar el estilo en el preambulo 
    % a [style=ieee]. Aqui lo mostramos como ejemplo visual.
    \begin{enumerate}
        \item[\small [1]] \small Manual de LaTeX, ``Símbolos Matemáticos'', [En línea]. Disponible en: \url{https://manualdelatex.com/simbolos}.
        \item[\small [2]] CodeCogs, ``Editor Visual de Ecuaciones'', 2026. [En línea]. Disponible en: \url{https://editor.codecogs.com}.
        \item[\small [3]] Tables Generator, ``Generador de tablas para LaTeX'', [En línea]. Disponible en: \url{https://www.tablesgenerator.com}.
    \end{enumerate}
\end{frame}

\begin{frame}{Referencias en Formato APA 7}
    \begin{itemize}
        \item[\quad] \small CodeCogs. (2026). \textit{Editor Visual de Ecuaciones}. Recuperado de \url{https://editor.codecogs.com}
        \item[\quad] \small Manual de LaTeX. (s.f.). \textit{Símbolos Matemáticos}. Recuperado de \url{https://manualdelatex.com/simbolos}
        \item[\quad] \small Tables Generator. (s.f.). \textit{Generador de tablas para LaTeX}. Recuperado de \url{https://www.tablesgenerator.com}
    \end{itemize}
\end{frame}

\end{document}
=======
\documentclass{beamer}

% Theme selection (you will need to choose or create a theme that matches your desired look)
%\usetheme{Madrid} % Example theme, you might want to explore others like CambridgeUS, Boadilla, etc.
% Or create a custom theme
% \usecolortheme{beaver} % Example color theme

\usepackage[utf8]{inputenc}
\usepackage{amsmath}
\usepackage{amsfonts}
\usepackage{amssymb}
\usepackage{graphicx} % Required for including images

% Custom commands or settings can go here to match the UNAD branding
% For example, setting specific colors:
% \definecolor{unadblue}{RGB}{X,Y,Z} % Replace X,Y,Z with actual RGB values
% \setbeamercolor{frametitle}{fg=unadblue}

\title[Transformada de Laplace]{WebConferencia Tarea 4 – Solución de Ecuaciones Diferenciales por Transformada de Laplace}
\author{Hugo López Aranda \\ Jonathan Alberto Cervantes}
\institute{Escuela de Ciencias Básicas, Tecnología e Ingeniería (ECBTI) \\ UNAD - Universidad Nacional Abierta y a Distancia}
\date{12 de junio de 2025}

\begin{document}

% Title slide
\begin{frame}
    \titlepage
    \begin{center}
        % You might want to include the UNAD logo here if you have it as an image file
        % \includegraphics[width=0.3\textwidth]{unad_logo.png}
    \end{center}
\end{frame}

% Table of Contents (optional)
\begin{frame}{Tabla de Contenidos}
    \tableofcontents
\end{frame}

% Introduction to Laplace Transform
\section{Transformada de Laplace}
\begin{frame}{Transformada de Laplace}
    \begin{itemize}
        \item Es una transformación matemática denotada por $\mathcal{L}$.
        \item Convierte funciones $f(t)$, dependientes del tiempo, en funciones $F(s)$, definidas en un dominio complejo.
        \item Convierte ecuaciones diferenciales en ecuaciones algebraicas, facilitando su solución.
        \item Aplicaciones en ingeniería: análisis y diseño de sistemas de control, transferencia de calor, y análisis de circuitos.
        \item Se define formalmente como: $\mathcal{L}[f(t)]=\int_{0}^{\infty}f(t)e^{-st}dt$.
    \end{itemize}
\end{frame}

% Inverse Laplace Transform
\begin{frame}{Transformada Inversa de Laplace}
    \begin{itemize}
        \item También conocida como Anti transformada.
        \item Toma funciones $F(s)$ y las retorna a funciones $f(t)$.
        \item Denotada por $\mathcal{L}^{-1}$.
    \end{itemize}
\end{frame}

% Example 1: Laplace Transform of e^(at)
\begin{frame}{Ejemplo 1: Transformada de Laplace}
    \centering
    \textbf{Si $f(t)=e^{at}$ entonces:}
    \[
    \mathcal{L}[f(t)]=\mathcal{L}[e^{at}]=\int_{0}^{\infty}(e^{at}*e^{-st})dt=\frac{1}{s-a}=F(s) \text{ }
    \]
    \textbf{En conclusión:} $\mathcal{L}[e^{at}]=\frac{1}{s-a}$ 
\end{frame}

% Example 2: Inverse Laplace Transform
\begin{frame}{Ejemplo 2: Transformada Inversa de Laplace}
    \centering
    \textbf{Si $F(s)=\frac{1}{s-a}$ entonces:}
    \[
    \mathcal{L}^{-1}[F(s)]=\mathcal{L}^{-1}\left[\frac{1}{s-a}\right]=e^{at}=f(t) \text{ }
    \]
    \textbf{En conclusión:} $\mathcal{L}^{-1}\left[\frac{1}{s-a}\right]=e^{at}$ 
\end{frame}

% Table of Transforms
\section{Tabla de Transformadas}
\begin{frame}{Tabla de Transformadas de Laplace}
    \centering
    Puede usar un recurso tecnológico para buscar más transformadas.

    \begin{columns}
        \column{0.5\textwidth}
        \textbf{Transformada de Laplace}
        \begin{itemize}
            \item $\mathcal{L}[kf(t)]=k\mathcal{L}[f(t)]$ (k constante) 
            \item $\mathcal{L}[f(t)\pm g(t)]=\mathcal{L}[f(t)]\pm\mathcal{L}[g(t)]$ (Linealidad) 
            \item $\mathcal{L}\{1\}=\frac{1}{s}$ 
            \item $\mathcal{L}\{e^{kt}\}=\frac{1}{s-k}$ 
            \item $\mathcal{L}\{t^{n}\}=\frac{n!}{s^{n+1}}$ 
            \item $\mathcal{L}\{t^{n}e^{kt}\}=\frac{n!}{(s-k)^{n+1}}$ 
            \item $\mathcal{L}[k]=\frac{k}{s}$ 
            \item $\mathcal{L}\{sen~kt\}=\frac{k}{s^{2}+k^{2}}$ 
            \item $\mathcal{L}\{cos~kt\}=\frac{s}{s^{2}+k^{2}}$ 
        \end{itemize}
        \column{0.5\textwidth}
        \textbf{Transformada de Laplace Inversa}
        \begin{itemize}
            \item $\mathcal{L}^{-1}\{\frac{1}{s}\}=1$ 
            \item $\mathcal{L}^{-1}\{\frac{1}{s-k}\}=e^{kt}$ 
            \item $\mathcal{L}^{-1}\{\frac{n!}{s^{n+1}}\}=t^{n}$ 
            \item $\mathcal{L}^{-1}\{\frac{n!}{(s-k)^{n+1}}\}=t^{n}e^{kt}$ 
            \item $\mathcal{L}^{-1}\{\frac{k}{s^{2}+k^{2}}\}=sen~kt$ 
            \item $\mathcal{L}^{-1}\{\frac{s}{s^{2}+k^{2}}\}=cos~kt$ 
        \end{itemize}
    \end{columns}
\end{frame}

% Example 3: Calculating Laplace Transform
\section{Ejemplos de Transformada de Laplace}
\begin{frame}{Ejemplo 3: Calcular $\mathcal{L}[f(t)]$ para $f(t)=5e^{2t}+e^{t}+2$}
    \begin{columns}
        \column{0.5\textwidth}
        \textbf{Proposición Enunciado o Expresión Matemática}
        \begin{itemize}
            \item $f(t)=5e^{2t}+e^{t}+2$ 
            \item $\mathcal{L}[f(t)]=\mathcal{L}[5e^{2t}+e^{t}+2]$ 
            \item $\mathcal{L}[5e^{2t}+e^{t}+2] = \mathcal{L}[5e^{2t}]+\mathcal{L}[e^{t}]+\mathcal{L}[2]$ 
            \item $\mathcal{L}[5e^{2t}]=5\mathcal{L}[e^{2t}]=5*\frac{1}{s-2}=\frac{5}{s-2}$ 
            \item $\mathcal{L}[e^{t}]=\frac{1}{s-1}$ 
            \item $\mathcal{L}[2]=\frac{2}{s}$ 
            \item $\mathcal{L}[f(t)]=\frac{5}{s-2}+\frac{1}{s-1}+\frac{2}{s}$ 
        \end{itemize}
        \column{0.5\textwidth}
        \textbf{Razón o Explicación}
        \begin{itemize}
            \item Función dada 
            \item Aplicación del símbolo de Transformada 
            \item Aplicación propiedad de linealidad 
            \item Aplicación de propiedades de la tabla 
            \item Aplicación de propiedades de la tabla 
            \item Aplicación de propiedades de la tabla 
            \item Consolidación 
        \end{itemize}
    \end{columns}
    \textbf{Conclusión:} $\mathcal{L}[f(t)]=\frac{5}{s-2}+\frac{1}{s-1}+\frac{2}{s}=F(s)$ 
\end{frame}

% Example 4: Calculating Laplace Transform with sin and cos
\begin{frame}{Ejemplo 4: Calcular $\mathcal{L}[f(t)]$ para $f(t)=t^{2}+3~sin(\pi t)+4t~cos(\pi t)$}
    \begin{columns}
        \column{0.5\textwidth}
        \textbf{Proposición Enunciado o Expresión Matemática}
        \begin{itemize}
            \item $f(t)=t^{2}+3~sin(\pi t)+4t~cos(\pi t)$ 
            \item $\mathcal{L}[f(t)]=\mathcal{L}[t^{2}+3~sin(\pi t)+4t~cos(\pi t)]$ 
            \item $\mathcal{L}[t^{2}+3~sin(\pi t)+4t~cos(\pi t)] = \mathcal{L}[t^{2}]+\mathcal{L}[3~sin~\pi t]+\mathcal{L}[4t~cos~\pi t]$ 
            \item $\mathcal{L}[t^{2}]=\frac{2!}{s^{2+1}}=\frac{2}{s^{3}}$ 
            \item $\mathcal{L}[3~sin~\pi t]=3\frac{\pi}{s^{2}+\pi^{2}}=\frac{3\pi}{s^{2}+\pi^{2}}$ 
            \item $\mathcal{L}[4t~cos~\pi t]=4\frac{s^{2}-\pi^{2}}{(s^{2}+\pi^{2})^{2}}=\frac{4(s^{2}-\pi^{2})}{(s^{2}+\pi^{2})^{2}}$ 
        \end{itemize}
        \column{0.5\textwidth}
        \textbf{Razón o Explicación}
        \begin{itemize}
            \item Función dada 
            \item Aplicación del símbolo de Transformada 
            \item Aplicación propiedad de linealidad 
            \item Aplicación de propiedades. 
            \item Aplicación de propiedades. 
            \item Aplicación de propiedades. 
        \end{itemize}
    \end{columns}
    \textbf{Conclusión:} $\mathcal{L}[f(t)]=\frac{2}{s^{3}}+\frac{3\pi}{s^{2}+\pi^{2}}+\frac{4(s^{2}-\pi^{2})}{(s^{2}+\pi^{2})^{2}}=F(s)$ 
\end{frame}

% Examples of Inverse Laplace Transform
\section{Ejemplos de Transformada Inversa de Laplace}
\begin{frame}{Transformada Inversa o Anti transformada de Laplace $\mathcal{L}^{-1}$}
    \begin{itemize}
        \item Otro concepto importante es la Transformada inversa de Laplace o Anti transformada que toma funciones $F(s)$ y las retorna a funciones $f(t)$.
    \end{itemize}
\end{frame}

% Example 5: Inverse Laplace Transform of k/(s^2+k^2)
\begin{frame}{Ejemplo 5: Calcular $\mathcal{L}^{-1}[F(s)]$ para $F(s)=[\frac{k}{s^{2}+k^{2}}]$}
    \centering
    \textbf{Proposición Enunciado o Expresión Matemática}
    \[
    \mathcal{L}^{-1}[F(s)]=\mathcal{L}^{-1}\left[\frac{k}{s^{2}+k^{2}}\right]=sin~kt \text{ }
    \]
    \textbf{Razón o Explicación}
    Lo anterior se da porque $\mathcal{L}[sin~kt]=\frac{k}{s^{2}+k^{2}}$ 
    \[
    \mathcal{L}^{-1}[\mathcal{L}[sin~kt]]=\mathcal{L}^{-1}\left[\frac{k}{s^{2}+k^{2}}\right] \text{ }
    \]
    \[
    sin~kt=\mathcal{L}^{-1}\left[\frac{k}{s^{2}+k^{2}}\right] \rightarrow \mathcal{L}^{-1}\left[\frac{k}{s^{2}+k^{2}}\right]=sin~kt \text{ }
    \]
    \textbf{Conclusión:} $f(t)=sin~kt$ 
\end{frame}

% Example 6: Inverse Laplace Transform of 2/(s+3)
\begin{frame}{Ejemplo 6: Calcular $\mathcal{L}^{-1}[\frac{2}{s+3}]$}
    \begin{columns}
        \column{0.5\textwidth}
        \textbf{Proposición Enunciado o Expresión Matemática}
        \begin{itemize}
            \item $F(s)=[\frac{2}{s+3}]$ 
            \item $\mathcal{L}^{-1}[\frac{2}{s+3}]$ 
            \item $2\mathcal{L}^{-1}[\frac{1}{s+3}]$ 
            \item $2\mathcal{L}^{-1}[\frac{1}{s+3}]=2e^{-3t}$ 
        \end{itemize}
        \column{0.5\textwidth}
        \textbf{Razón o Explicación}
        \begin{itemize}
            \item Función dada 
            \item Aplicación del símbolo de Transformada inversa. 
            \item Aplicación de Linealidad para la Transformada inversa 
            \item Se busca en la tabla de transformadas una expresión parecida a $\mathcal{L}^{-1}[\frac{1}{s+3}]$, la cual es $\mathcal{L}[e^{at}]=\frac{1}{s-a}=\frac{1}{s-(-3)}=\frac{1}{s+3}$ donde $a=(-3)$ 
        \end{itemize}
    \end{columns}
    \textbf{Conclusión:} $f(t)=2e^{-3t}$ 
\end{frame}

% Example 7: Inverse Laplace Transform of 1/s^5
\begin{frame}{Ejemplo 7: Calcular $\mathcal{L}^{-1}[\frac{1}{s^{5}}]$}
    \begin{columns}
        \column{0.5\textwidth}
        \textbf{Proposición Enunciado o Expresión Matemática}
        \begin{itemize}
            \item $F(s)=[\frac{1}{s^{5}}]$ 
            \item $\mathcal{L}^{-1}[\frac{1}{s^{5}}]$ 
            \item $\mathcal{L}^{-1}[\frac{1}{s^{(4+1)}}]$ 
            \item $\mathcal{L}^{-1}[\frac{1}{s^{(4+1)}}]=\mathcal{L}^{-1}[\frac{4!}{s^{(4+1)}}*\frac{1}{4!}]$ 
            \item $\frac{1}{4!}\mathcal{L}^{-1}[\frac{4!}{s^{(4+1)}}]$ 
            \item $\frac{1}{24}t^{4}$ 
        \end{itemize}
        \column{0.5\textwidth}
        \textbf{Razón o Explicación}
        \begin{itemize}
            \item Función dada 
            \item Aplicación del símbolo de Transformada inversa. 
            \item Se busca en la tabla de transformadas y anti transformadas una expresión que se parezca a $\mathcal{L}^{-1}[\frac{1}{s^{(4+1)}}]$, la cual es $\mathcal{L}[t^{n}]=\frac{n!}{s^{n+1}}=\frac{4!}{s^{4+1}}$ donde $n=4$ 
            \item Multiplicando y a su vez dividiendo entre 4! 
            \item Aplicando linealidad para $\mathcal{L}^{-1}$ 
        \end{itemize}
    \end{columns}
    \textbf{Conclusión:} $f(t)=\frac{1}{24}t^{4}$ 
\end{frame}

% Partial Fractions
\section{Fracciones Parciales}
\begin{frame}{Fracciones Parciales}
    \begin{itemize}
        \item Las fracciones parciales son una técnica algebraica que se utiliza para descomponer una fracción racional compleja en una suma de fracciones más simples.
        \item Esto facilita integraciones y transformadas inversas como la de Laplace.
    \end{itemize}
\end{frame}

% Example 8: Partial Fraction Decomposition
\begin{frame}{Ejemplo 8: Descomponer en fracciones parciales la expresión $\frac{1}{s(s+2)}$}
    \centering
    Es válido usar un recurso tecnológico para hacer la descomposición. 
    \textbf{Resultado final:}
    \[
    \frac{1}{s(s+2)}=\frac{1}{2s}-\frac{1}{2(s+2)} \text{ }
    \]
\end{frame}

% Example 9: Application of Partial Fractions in Inverse Laplace Transform
\begin{frame}{Ejemplo 9: Aplicación de las fracciones parciales en la Transformada inversa de Laplace.}
    \centering
    \textbf{Hallar la transformada inversa de Laplace para la siguiente función:} $F(s)=\frac{1}{s(s+2)}$ 

    \begin{columns}
        \column{0.5\textwidth}
        \textbf{Proposición Enunciado o Expresión Matemática}
        \begin{itemize}
            \item $F(s)=\frac{1}{s(s+2)}$
            \item $\frac{1}{s(s+2)} = \frac{1}{2s}-\frac{1}{2(s+2)}$
            \item $\mathcal{L}^{-1}\left[\frac{1}{2s}-\frac{1}{2(s+2)}\right]$
            \item $\frac{1}{2}\mathcal{L}^{-1}\left[\frac{1}{s}\right]-\frac{1}{2}\mathcal{L}^{-1}\left[\frac{1}{s+2}\right]$
            \item $\frac{1}{2}(1) - \frac{1}{2}e^{-2t}$
        \end{itemize}
        \column{0.5\textwidth}
        \textbf{Razón o Explicación}
        \begin{itemize}
            \item Enunciado a resolver 
            \item Aplicación de fracciones parciales con recurso tecnológico 
            \item Aplicación del símbolo de Transformada inversa sobre las fracciones parciales 
            \item Aplicación de linealidad 
            \item De la tabla de transformadas inversas sabemos que $\mathcal{L}^{-1}[\frac{1}{s}]=1$ y $\mathcal{L}^{-1}[\frac{1}{s+2}]=e^{-2t}$ 
        \end{itemize}
    \end{columns}
    \textbf{Conclusión:} $f(t)=\frac{1}{2}-\frac{1}{2}e^{-2t}$ 
\end{frame}

% Suggestion for Exercise 5
\section{Sugerencias}
\begin{frame}{Sugerencia para el Ejercicio 5 Tarea 4}
    \textbf{Anotación:} Para el desarrollo del Ejercicio 5 de su Tarea 4 puede trabajar con los videos propuestos.
    \textbf{Videos Propuestos:}
    \begin{itemize}
        \item \url{https://www.youtube.com/watch?v=t8JdvZC63bY} 
        \item \url{https://www.youtube.com/watch?v=tv_5Khz4GJs} 
        \item \url{https://www.youtube.com/watch?v=LpQ5EkWLsDc} 
    \end{itemize}
    O buscar algún video referido al tema de aplicaciones matemáticas en el canal de la UNAD: \url{https://www.youtube.com/universidadunad} 
\end{frame}

% Thank You slide
\begin{frame}
    \centering
    \Huge ¡GRACIAS! 
    \vspace{1cm}
    \Large \url{www.unad.edu.co} 
    \vspace{0.5cm}
    % \includegraphics[width=0.4\textwidth]{unad_social_media.png} % Replace with actual social media icons if you have them as an image
\end{frame}

\end{document}
>>>>>>> makefile
